\documentclass[10pt,a4paper]{amsart}
\usepackage[margin=19mm]{geometry}
\usepackage{amsmath,amssymb,mathtools}
\usepackage[scr=rsfs,cal=euler]{mathalfa}
\usepackage{xfrac}
\usepackage[unicode,hidelinks,bookmarks=false]{hyperref}
\newcommand{\ie}{i.e.\ }
\newcommand{\cf}{cf.\ }
\newcommand{\resp}{respectively\ }
\newcommand{\Subsection}{\subsection}
\providecommand{\nobreakdash}{\nobreak\hbox{-}}
\setlength{\emergencystretch}{2em}
\multlinegap0pt
\newcommand{\R}{\ensuremath{\mathbb{R}}}
\newcommand{\Z}{\ensuremath{\mathbb{Z}}}
\newcommand{\G}{\ensuremath{\mathbb{H}}}
\newcommand{\N}{\ensuremath{\mathbb{N}}}
\newcommand{\K}{\ensuremath{\mathbb{K}}}
\newcommand{\PT}{\ensuremath{\mathbb{P}}}
\newcommand{\T}{\ensuremath{\mathbb{T}}}
\newcommand{\C}{\ensuremath{\mathbb{C}}}
\newcommand{\I}{\ensuremath{\mathbb{I}}}
\newcommand{\A}{\ensuremath{\mathbb{A}}}
\newcommand{\X}{\ensuremath{\mathbb{X}}}
\newcommand{\D}{\ensuremath{\mathbb{D}}}
\newcommand{\floor}[1]{\lfloor #1 \rfloor}
\newcommand{\Cerg}{\ensuremath{\mathbb{C}_{\textup{erg}}}}
\newcommand{\Chilg}{\ensuremath{\mathbb{C}_{\mathcal H}}}
\newcommand{\Rerg}{\ensuremath{\gamma}}
\newcommand{\Ierg}{\ensuremath{\omega}}
\DeclareMathOperator{\myRe}{Re}
\DeclareMathOperator{\myIm}{Im}
\numberwithin{equation}{section}
\newtheorem{theorem}{Theorem}[section]
\newtheorem{lemma}[theorem]{Lemma}
\newtheorem{definition}[theorem]{Definition}
\newtheorem{proposition}[theorem]{Proposition}
\newtheorem{remark}[theorem]{Remark}
\newtheorem{example}[theorem]{Example}
\allowdisplaybreaks
\title[Ergodic complex plane and cylinder transformation: original Theorem 3.5]{Ergodic complex plane and cylinder transformation on a periodic time scale: original Theorem 3.5 and its complete original proof}
\author{John M. Davis, Billy J. Jackson and Dylan Poulsen}
\date{Source: Electronic Journal of Differential Equations 2026 (2026), No. 04, pp. 1--20; DOI 10.58997/ejde.2026.04}
\begin{document}
\maketitle
\noindent\textit{Editorial scope.} The counted unit of this excerpt is the article's own Theorem 3.5 as printed, together with its complete original proof. Everything the proof invokes is retained verbatim as uncounted context and is not counted again: the whole of the original Section 1 (Hilger's complex plane, the cylinder transformation, the regressive points, condition (C1) and the set $\mathcal S(\mathbb T)$ of exponential stability, the periodic time scales $\mathbb T_{\mu_1,\dots,\mu_n}$ and $\mathbb P_{a,b}$), the whole of the original Section 2 (the ergodic growth rate $\gamma_{B^c}$ and ergodic frequency $\omega_{B^c}$ and the standing hypothesis (P)), and the whole of the original Section 3 up to the counted theorem, which contains Definition 3.1, Theorems 3.2--3.3 with the complete proof of Theorem 3.3, and Lemma 3.4. The numbering of the retained statements is the published one, so the printed numbers agree with the publisher PDF. One illustration lies inside the excerpt: Figure 1 of the article. Its caption, its label and the cross-reference to it are kept verbatim, but the external asset file is not redistributed and its \verb|\includegraphics| line is the single omitted source string, recorded under a bounded-source-comparison fidelity receipt. Printed slips are kept literal and no mathematics has been rewritten or repaired.
\section{Hilger's complex plane and exponential stability}

We begin with a brief summary of the Hilger complex plane \cite{BoPeD} on a time
scale $\mathbb{T}$. To match the narrative of later sections, we frame the Hilger
complex plane entirely in terms of the cylinder transformation, which is a different
perspective from \cite{BoPeD}.

Let $\mathbb{T}$ be a time scale and let $t \in \mathbb{T}$. Define the \emph{Hilger complex numbers}
at $t\in\mathbb{T}$, denoted $\mathbb{C}_{\mu(t)}$, by $\mathbb{C}_{\mu(t)} = \mathbb{C} \setminus \{-1/\mu(t)\}$. Defining $\Omega(\mu(t)) := \pi/\mu(t)$, let the \emph{cylinder strip} at $t\in\mathbb{T}$ be defined as
\[
\mathbb{C}^{\Omega(\mu(t))} := \{z \in \mathbb{C} \mid -\Omega(\mu(t)) < \myIm(z) \leq \Omega(\mu(t)) \}.
\]
\begin{definition} \rm
Let $\mathbb{T}$ be a time scale and let $t \in \mathbb{T}$. Define the
\emph{cylinder transformation} $\xi_{\mu(t)}:\mathbb{C}_{\mu(t)} \to \mathbb{C}^{\Omega(\mu(t))}$ by
 \begin{equation}\label{cyl}
 \xi_{\mu(t)}(z):=
 \begin{cases}
 \frac{\operatorname{Log}(1+z \mu(t))}{\mu(t)}, & \mu(t)>0,\\
 z, & \mu(t)=0.
 \end{cases}
 \end{equation}
\end{definition}

Note that $\xi_{\mu(t)}$ is one-to-one, and so $\xi_{\mu(t)}^{-1}$ exists.
In fact, it has the explicit form
\begin{equation}\label{invcyl}
 \xi_{\mu(t)}^{-1}(z):=
 \begin{cases}
 \frac{e^{\mu(t) z} -1}{\mu(t)}, & \mu(t)>0,\\
 z, & \mu(t)=0.
 \end{cases}
 \end{equation}

The \emph{Hilger real part} at $t\in\mathbb{T}$ of $z \in \mathbb{C}_{\mu(t)}$ can be defined as
 $$
 \myRe_{\mu(t)}(z)= \xi_{\mu(t)}^{-1}(\myRe(\xi_{\mu(t)}(z))).
 $$
The \emph{Hilger imaginary part} at $t \in \mathbb{T}$ of $z$ is given by
 $$
 \myIm_{\mu(t)}(z) = \myIm(\xi_{\mu(t)}(z)),
 $$
while the \emph{Hilger pure imaginary part} of $z$ at $t \in \mathbb{T}$ is given by
 $$
 \stackrel{o}{\iota} \myIm_{\mu(t)}(z) := \xi^{-1}_{\mu(t)}( i \myIm(\xi_{\mu(t)}(z))).
 $$
It is straightforward to establish the following identities (which are usually taken as definitions).

\begin{proposition} \label{lem:HilXi}
Let $\mathbb{T}$ be a time scale and let $t \in \mathbb{T}$. Then
\begin{gather*}
\myRe_{\mu(t)}(z)  = \lim_{\tau \searrow \mu(t)} \frac{|1+\tau z|-1}{\tau}, \\
\myIm_{\mu(t)}(z)  = \lim_{\tau \searrow \mu(t)}\frac{\operatorname{Arg}(z\tau+1)}{\tau}, \\
\stackrel{o}{\iota} \myIm_{\mu(t)}(z)
 = \lim_{\tau \searrow \mu(t)}\frac{e^{i \myIm_{\mu(t)}(z) \tau}-1}{\tau},
\end{gather*}
where $\operatorname{Arg}$ is the principal argument satisfying $-\pi < \operatorname{Arg}(z) \leq \pi$.
\end{proposition}

The \emph{circle plus operation} $\oplus_{\mu(t)}$ is defined by
$a \oplus_{\mu(t)} b:=a+b+\mu(t) ab$. The \emph{circle minus operation}
$\ominus_{\mu(t)}$ is defined to be the additive inverse of $\oplus_{\mu(t)}$.
Bohner and Peterson
\cite[Theorem~2.24]{BoPeD}  establish that $(\mathbb{C}_{\mu(t)}, \oplus_{\mu(t)})$
is homomorphic to $(\mathbb{C}^{\Omega(\mu(t))},+ (\text{mod } 2 \pi i/\mu(t)))$ with
group homomorphism $\xi_{\mu(t)}$. Although not explicitly stated in \cite{BoPeD},
it follows from the fact that
$\xi_{\mu(t)}$ is a bijection that
$(\mathbb{C}_{\mu(t)}, \oplus_{\mu(t)})$ is isomorphic to
$(\mathbb{C}^{\Omega(\mu(t))},+ (\text{mod } 2 \pi i/\mu) )$.

The group isomorphism $\xi_{\mu(t)}$ implies we can understand the circle plus
addition via, for $w,z \in \mathbb{C}_{\mu(t)}$,
\begin{equation}
\label{eq:oplusViaXi}
z \oplus_{\mu(t)} w = \xi_{\mu(t)}^{-1}(\xi_{\mu(t)}(z)
+ \xi_{\mu(t)}(w) \pmod{2 \pi i/\mu(t)}).
\end{equation}

Using circle plus, for fixed $t \in \mathbb{T}$, every $z \neq -\frac{1}{\mu(t)}$ has the
\emph{Hilger decomposition at time $t$} given by
$$
z=\text{Re}_{\mu(t)}(z) \oplus_{\mu(t)} \stackrel{o}{\iota}
\operatorname{Im}_{\mu(t)}(z).
$$

The Hilger complex plane, $\mathbb{C}_{\mu(t)}$, is shown in
Figure~\ref{fig:complexplane} for $\mu(t)\neq 0$.
Important components of $\mathbb{C}_{\mu(t)}$ are the \emph{Hilger disk}
at $t\in \mathbb{T}$, given by
$$
 \mathcal{H}_{\mu(t)} := \{ z \in \mathbb{C}_{\mu(t)} \mid \myRe_{\mu(t)}(z) < 0 \},
$$
which can be thought of as the preimage of the left half-plane of
$\mathbb{C}^{\Omega(\mu(t))}$ under $\xi_{\mu(t)}$.
The \emph{Hilger circle} at $t \in \mathbb{T}$ is the boundary of $\mathcal{H}_{\mu(t)}$,
which we denote by $\mathbb{I}_{\mu(t)}$.
The Hilger circle can be thought of as the preimage of the imaginary axis of
$\mathbb{C}^{\Omega(\mu(t))}$ under $\xi_{\mu(t)}$.


\begin{figure}[t]
  \centering
  
  \caption{Hilger's complex plane, $\mathbb{C}_{\mathcal{H}}$. The $z$ inside the circle have negative
  Hilger real part, the $z$ on the circle have zero Hilger real part, and the
  $z$ outside the circle have positive Hilger real part. Points $z$ on the
  Hilger real axis $\mathbb{R}_{\mu(t)}$ (the solid ray on the real axis) are such
  that $e_{z}(t,t_0)>0$ for all $t$, while points on the Hilger alternating axis
  $\mathbb{A}_{\mu(t)}$ (the dotted ray on the real axis) are such that $e_{z}(t,t_0)$
  changes sign at each $t \in \mathbb{T}$.}
  \label{fig:complexplane}
\end{figure}



The cylinder transformation maps contours of equal Hilger imaginary part
(which are rays emanating from $-1/\mu(t)$) to horizontal lines in
$\mathbb{C}^{\Omega(\mu(t))}$. Also, contours of equal Hilger real part
(which are circles centered at $-1/\mu(t)$) are mapped to vertical lines in
$\mathbb{C}^{\Omega(\mu(t))}$. In particular, these contours form an orthogonal curvilinear coordinate system on $\mathbb{C}_{\mu(t)}$.

\begin{remark} \rm
The cylinder transformation is strongly related to the concepts of growth rate and
frequency on time scales with constant graininess. On the time scale
$\mathbb{T}=h\mathbb{Z}$, the contours of equal Hilger imaginary part correspond
to frequency given by
\begin{equation}
\label{eq:usefulIm}
\omega_h(z):=\operatorname{Im}_h(z) = \myIm(\xi_h(z)) = \frac{\operatorname{Arg}(1+z h)}{h}.
\end{equation}
This is the same quantity that represents frequency in Hilger's definition of the
time scale sinusoids \cite{HiA}. Moreover, the contours of equal Hilger real part
all have the same exponential growth rate corresponding to the Lyapunov exponent
\begin{equation}
\label{eq:usefulRe}
\gamma_{h}(z):= \frac{\ln(1+h\, \text{Re}_h(z))}{h}
= \xi_{h}(\myRe_h(z)) = \myRe(\xi_h(z)) = \frac{\ln|1+hz|}{h}.
\end{equation}
\end{remark}
Similarly, on the time scale $\mathbb{T}=\mathbb{R}$, $\xi_0(z) = z$, and the contours
of equal exponential growth rate are vertical lines determined by the Lyapunov
exponent $\gamma_0(z):=\text{Re}(z)$, while the contours of equal frequency are
horizontal lines given by $\omega_0(z)=\operatorname{Im}(z)$.

We say that $p:\mathbb{T} \to \mathbb{R}$ is \emph{regressive} if $1+\mu(t)p(t)\not=0$
for all $t\in\mathbb{T}$.
For regressive $p$, the \emph{time scale exponential function} $e_p(t,s)$ is defined by
 \[
 e_{p}(t,s) := \exp\Big(\int_{s}^{t} \xi_{\mu(t)}(p(\tau)) \Delta \tau \Big).
 \]
Following Karpuz \cite{Ka}, denote the set of all \emph{regressive points} in the
complex plane by
$$
\mathbb{C}_{\mu(\mathbb{T})}:=\{z \in \mathbb{C} \mid 1+\mu(t) z \neq 0 \text{ for all } t \in \mathbb{T}\}.
$$

The definitions related to Hilger's complex plane are \emph{local} and \emph{dynamic}
in nature; that is, they depend on and change in $t$. In 2003, P\"{o}tzsche, Siegmund,
 and Wirth introduced the concept of a \emph{global} and \emph{static} exponential
 growth rate for the time scale exponential function.

\begin{theorem}[P\"{o}tzsche, Siegmund, and Wirth \cite{PoSiWi}]\label{PSW1}
Let $\mathbb{T}$ be a time scale which is unbounded above and let $\lambda
\in \mathbb{C}$. Then the scalar equation
 $$
 x^{\Delta}(t)=\lambda x(t), \quad x(t_0)=x_0,
 $$
is exponentially stable if and
only if one of the following conditions is satisfied for arbitrary
$t_0 \in \mathbb{T}$:
\begin{itemize}
\item[(C1)] $\gamma(\lambda):=  \limsup_{T \to
\infty}\frac{1}{T-t_0}\int_{t_0}^{T}\lim_{s \searrow
\mu(t)} \frac{\ln|1+\lambda s|}{s}\Delta t <0$,

\item[(C2)] For every
$T \in \mathbb{T}$, there exists a $t \in \mathbb{T}$ with $t>T$ such that
$1+\mu(t)\lambda=0$,
\end{itemize}
where we use the convention $\ln 0=-\infty$ in \textup{(C1)}.
\end{theorem}

Note that condition (C1) can be restated using \eqref{eq:usefulRe} as
\begin{itemize}
\item[(C1*)] $ \limsup_{T \to
\infty}\frac{1}{T-t_0}\int_{t_0}^{T}\lim_{s \searrow
\mu(t)} \gamma_s(\lambda) \Delta t <0$.
\end{itemize}
This theorem naturally leads to the following definition.

\begin{definition}[\cite{PoSiWi}] \rm
Given a time scale $\mathbb{T}$ which is unbounded above, for
arbitrary $t_0 \in \mathbb{T}$, define the sets
 \begin{gather*}
 \mathcal{S}_{\mathbb{C}}(\mathbb{T}):=\{\lambda \in \mathbb{C}: \text{ (C1) holds}\},\\
 \mathcal{S}_{\mathbb{R}}(\mathbb{T}):=\{\lambda \in \mathbb{R}: \text{ (C2) holds}\}.
 \end{gather*}
Then the \emph{set of exponential stability} for $\mathbb{T}$ is given by
$$
\mathcal{S}(\mathbb{T}):=\mathcal{S}_{\mathbb{C}}(\mathbb{T}) \cup \mathcal{S}_{\mathbb{R}}(\mathbb{T}).
$$
\end{definition}

The calculation of the set of exponential stability is simplified for periodic
time scales [\cite[Lemma~3.4(b)]{PoSiWi}, , which will be a focus of this paper.

\begin{definition} \rm
 A time scale $\mathbb{T}$ is \emph{periodic} if there exists $L>0$ such that
 $t \in \mathbb{T}$ implies $t+L \in \mathbb{T}$ and $\mu(t+L)=\mu(t)$ for all $t \in \mathbb{T}$.
We call $L$ a \emph{period} of the time scale.
\end{definition}


\begin{remark} \rm
Throughout this paper, we will illustrate various results on the following
prototypical time scales.
 \begin{enumerate}
 \item For $\mu_1,\mu_2,\dots,\mu_n>0$, set $L=\mu_1+\mu_2+\cdots+\mu_n$.
 Then $\mathbb{T}_{\mu_1,\mu_2,\dots,\mu_n}$ is the time scale starting at 0 whose
 graininesses form the periodic sequence $\{\mu_1,\mu_2,\dots,\mu_n\}$, i.e.,
 $$
 \mathbb{T}_{\mu_1,\mu_2,\dots,\mu_n}
 =\Big\{0,\mu_1,\mu_1+\mu_2,\dots,\underbrace{\mu_1+\cdots+\mu_n}_{L},
 L+\mu_1,L+\mu_1+\mu_2,\dots,2L,\dots\Big\}.
 $$

 \item $\mathbb{P}_{a,b}$ is the periodic time scale starting at 0 with an
 interval of length $a$ followed by a gap of length $b$, i.e.,
 $$
 \mathbb{P}_{a,b}=[0,a]\cup[a+b,2a+b]\cup[2a+2b,3a+2b]\cup\cdots
 $$
\end{enumerate}
\end{remark}

A key insight for this paper is that since the cylinder transformation is linked to both
growth rate (via the Hilger real part) and frequency (via the Hilger imaginary part) on
constant graininess time scales, and since the average of the real part of the cylinder
transformation is related to global exponential growth rate, it seems reasonable that
the average of the imaginary part of the cylinder transformation is related to the
global frequency. 


\section{Ergodic growth rate and frequency}

To solidify the connection between the average of the cylinder transformation and the
growth rate and frequency of the time scale exponential, notice for $\lambda \in \mathbb{C}_h$
that the time scale exponential function decomposes as
\begin{equation} \label{eq:exp_gamma_omega}
\begin{aligned}
 e_{\lambda}(t,t_0)
 &=  \exp\Big(\int_{t_0}^{t} \xi_{\mu(\tau)}(\lambda) \Delta \tau \Big) \\
 &= \exp\Big( \frac{\int_{t_0}^{t} \xi_{\mu(\tau)}(\lambda) \Delta \tau}
  {t-t_0} (t-t_0) \Big) \\
 &= \exp\Big( \frac{\int_{t_0}^{t} \myRe(\xi_{\mu(\tau)}(\lambda))
  \Delta \tau}{t-t_0} (t-t_0) \Big)
  \exp\Big( i \frac{\int_{t_0}^{t} \myIm(\xi_{\mu(\tau)}(\lambda)) \Delta \tau}{t-t_0}
   (t-t_0) \Big) \\
 &\overset{\eqref{eq:usefulIm},\eqref{eq:usefulRe}}{=}
 \exp\Big( \frac{\int_{t_0}^{t} \gamma_{\mu(\tau)}(\lambda) \Delta \tau}{t-t_0}
  (t-t_0) \Big)
  \exp\Big( i \frac{\int_{t_0}^{t} \omega_{\mu(\tau)}(\lambda) \Delta \tau}{t-t_0}
  (t-t_0) \Big) \\
 &=  \exp\Big( \frac{\int_{t_0}^{t} \gamma_{\mu(\tau)}(\lambda) \Delta \tau}{t-t_0}
  (t-t_0) \Big)   \\
 &\quad \times \Big[\cos\Big( \frac{\int_{t_0}^{t} \omega_{\mu(\tau)}(\lambda)
 \Delta \tau}{t-t_0} (t-t_0) \Big) + i \sin\Big( \frac{\int_{t_0}^{t}
 \omega_{\mu(\tau)}(\lambda) \Delta \tau}{t-t_0} (t-t_0) \Big) \Big].
 \end{aligned}
\end{equation}
Therefore, over an interval $[t,t_0]_{\mathbb{T}}$, the growth rate and frequency of the
time scale exponential function are averages (in the integral sense) of
$\gamma_{\mu(\tau)}$ and $\omega_{\mu(\tau)}$, respectively.

Since we are averaging the cylinder transformation, which is a logarithm for
$\mu(t)>0$, we want to be careful about the branch cut. We begin by defining a subset
of the complex plane which avoids the branch cuts of \eqref{cyl} for all $\tau\in\mathbb{T}$.
Assuming \eqref{P}, $\mu_{\max}:=\max \{ \mu(t) \mid t \in \mathbb{T} \}$ exists, so define
\[
 B:= (-\infty, -1/\mu_{\max}],\quad
 B^c= \mathbb{C} \setminus B.
\]
Motivated by the form \eqref{eq:exp_gamma_omega}, we now formalize the averaged
quantities introduced above.

\begin{definition}
\label{def_gamma_omega} \rm
Let $\mathbb{T}$ be a time scale and $x,y \in \mathbb{R}$ such that $x+iy \in B^c$.
The \emph{ergodic growth rate of $e_{x+iy}(t,t_0)$ on $[t_0,t]_\mathbb{T}$ off the branch cut}
is given by
\begin{equation}
 \label{eq:gammatt0}
 \gamma_{B^c}(x,y,t_0,t) :=\frac{1}{t-t_0}\int_{t_0}^{t} \gamma_{\mu(\tau)}(x+iy)
 \Delta \tau
\end{equation}
Similarly, the \emph{ergodic frequency of $e_{x+iy}(t,t_0)$ on $[t_0,t]_\mathbb{T}$ off the
branch cut} is given by
 \begin{equation}
 \omega_{B^c}(x,y,t_0,t) := \frac{1}{t-t_0}\int_{t_0}^{t}
 \omega_{\mu(\tau)}(x+iy)\,\Delta \tau.\label{w}
 \end{equation}
\end{definition}

\begin{remark} \rm
Several remarks are in order regarding $\gamma_{B^c}(x,y,t_0,t)$ and
$\omega_{B^c}(x,y,t_0,t)$.

(1) We use the term \emph{ergodic} to describe the growth rate and frequency here
because, if one considers the space $\mathbb{X}=\mu^*(\mathbb{T})$ of all extended graininesses
of $\mathbb{T}$ \cite{JaDaE} in the order that they appear in $\mathbb{T}$, then the global (spatial)
averages of the growth rates and frequencies derived from $\mathbb{X}$ are the
pointwise limit of the local (time) averages of the local growth rates and frequencies
derived from $\mathbb{X}$ as defined above. We may treat the collection of asymptotic
extended graininesses as an irreducible Markov chain, which guarantees the ergodicity
of the associated stochastic process. See \cite{PoDaGr}.

(2) By \eqref{eq:usefulRe}, $\gamma_{h}(x+iy)$ has a few useful forms. The choice of
form gives different interpretations to the equation. Writing
$\gamma_{h}(x+iy) = \xi_{h}(\myRe_{h}(x+iy))$ is useful when a characterization of
growth rates in terms of the Hilger real part of $x+iy$ is advantageous.
By contrast, writing $\gamma_{h}(x+iy) = \myRe(\xi_{h}(x+iy))$ is in line with
P\"otzsche, Siegmund, and Wirth's condition.

(3) Using properties of the logarithm, we get
$$
\gamma_{h}(x+iy) = \ln(1+h \myRe_{h}(x+iy))^{1/h}.
$$
This implies $\gamma_{B^c}(x,y,t_0,t)$ is the natural log of the geometric mean (in
the time scales sense) on $[t_0,t]_\mathbb{T}$ of
\begin{equation}\label{loc}
 \lim_{s \searrow\mu(t)}(1+\operatorname{Re}_{s}(\lambda) s)^{1/s}
 = \lim_{s \searrow\mu(t)}|1+\lambda s|^{1/s}.
\end{equation}
Thus, \eqref{loc} can be interpreted as a type of local growth rate of $e_{x+iy}(t,t_0)$
at time $t$. Since the geometric mean is the best measure of average local growth rates,
expressing $\gamma_{h}(x+iy)$ in this way shows that \eqref{eq:gammatt0} is an effective
 measure of the average growth rate of the time scale exponential function over
$[t_0,t]_{\mathbb{T}}$. This observation was made earlier in \cite{DaGrMaJa,JaDaE} --albeit
from the geometric perspective-- by geometrically averaging Hilger circles.

(4) The formulation in \eqref{eq:exp_gamma_omega} is similar to the one provided by
 Karpuz \cite{Ka}:
 \[% \label{karpuz}
 e_{z}(t,t_0) = e_{\text{Re}_{\mu(t)}(z)}(t,t_0)
 \Big[ \cos\Big(\int_{t_0}^{t} \operatorname{Im}_{\mu(\eta)}(z) \Delta \eta\Big)
 + i \sin\Big(\int_{t_0}^{t} \operatorname{Im}_{\mu(\eta)}(z) \Delta \eta\Big) \Big].
 \]
To use the terminology of \cite{JaDalap}, Karpuz gave a type II (time scale)
exponential representation whereas \eqref{eq:exp_gamma_omega} is a type I (continuous)
exponential representation.

(5) In this new notation, condition (C1*) becomes
 \begin{equation} \label{supgamma}
 \limsup_{T \to \infty} \gamma_{B^c}(x,y,t_0,T) < 0,
 \end{equation}
and thus (C1*) is indeed a requirement that the largest cluster point of the average
exponential growth rate over the tail of the time scale is negative.
Therefore, \eqref{supgamma} is a condition on the \emph{global} exponential growth rate.
 Moreover, the formulation here ties $\gamma_{B^c}(x,y,t_0,t)$ to the local Lyapunov
 exponent $\gamma_{\mu(t)}$, in the spirit of P\"otzsche, Siegmund, and Wirth's
 (C1) condition.

(6) Note that the dynamic equation
\begin{equation}
z^{\Delta}(t)=\lambda z(t), \quad z(t_0)=z_0,
\label{expequation}
\end{equation}
for $\lambda \in B^c$ has a growth rate of $\gamma_{B^c}(x,y,t_0,\infty)$ rather than a growth rate of $\lambda$. Furthermore, unlike on $\mathbb{R}$ where the terms \emph{growth rate} and \emph{exponential order} are often used synonymously, these two quantities will be distinct in general. Indeed, if $\lambda$ in \eqref{expequation} is real and positive, then $\lambda$ is best described as the \emph{exponential order} of the equation.
 \end{remark}

We now see that we can understand the \emph{global} growth rate and frequency of the
time scale exponential function by understanding the asymptotic behavior of the
ergodic growth rate and ergodic frequency as $t\to\infty$.
 Since this asymptotic behavior is complicated in general (indeed, the limit as
$t\to\infty$ may not exist), we begin by studying the asymptotic behavior with the
following simplifying assumption:
\begin{equation}
\mathbb{T} \text{ is a periodic time scale with a period of } L>0
\text{ and }0 \in \mathbb{T}.
\tag{P}
\label{P}
\end{equation}
Under this assumption, we have
\[
\lim_{t \to \infty} \gamma_{B^c}(x,y,t,0) = \gamma_{B^c}(x,y,L,0),
\]
and similarly
\[
\lim_{t \to \infty} \omega_{B^c}(x,y,t,0) = \omega_{B^c}(x,y,L,0).
\]
Therefore, we can simplify notation by rewriting \eqref{eq:gammatt0}, \eqref{w},
for $x+iy \in B^c$, as
\begin{gather}
 \gamma_{B^c}(x,y) :=\frac{1}{L}\int_0^{L} \lim_{s \searrow
\mu(\tau)} \gamma_s(x+iy)\,\Delta \tau, \label{g21} \\
\omega_{B^c}(x,y) := \frac{1}{L}\int_{0}^{L} \lim_{s \searrow
\mu(\tau)} \omega_s(x+iy)\,\Delta \tau.\label{w1}
\end{gather}

\begin{remark} \rm
While the assumption of periodicity may appear to be limiting, a larger class of
aperiodic time scales can be analyzed using periodic techniques.
We will call a time scale \emph{simple} if there is a representative periodic time
scale which has the same limiting functions $\gamma_{B^c}$ and $\omega_{B^c}$.

To illustrate this, consider the symmetric time scale
$\mathbb{T}_{\text{PTM}} =\{0,t_1,t_2,\ldots\}$ with the graininess function $\mu(t_n)$
 equal to the $n^{\text{th}}$ term of the celebrated Prouhet-Thue-Morse sequence on
the symbols one and two \cite{allouche1999ubiquitous}.
Since the Prouhet-Thue-Morse sequence is aperiodic, $\mathbb{T}_{\text{PTM}}$ is also aperiodic.
Additionally, since the symbols of the sequence appear with equal weight in the limit
as the sequence length increases,
\begin{align*}
 \limsup_{T \to \infty} \gamma_{B^c}(x,y,t_0,T)
&= \limsup_{T \to \infty} \frac{1}{T-t_0} \int_{t_0}^{T} \frac{\ln|1+(x+iy) s|}{s} \Delta \tau\\
&= \frac{1}{3} \left[ \ln|1+(x+iy)| + \ln|1+2(x+iy)| \right].
\end{align*}
This last expression is equal to $\gamma_{B^c}(x,y)$ on the periodic time scale
$\mathbb{T}_{1,2}$. Similar results hold for calculating the limiting value of $\omega_{B^c}$.

In this paper, any result which assumes \eqref{P} can be extended to a simple time
scale via its representative (asymptotically equivalent) periodic time scale since our
arguments rely only on the limiting values of $\gamma_{B^c}$ and $\omega_{B^c}$
rather than the precise nature of $\mathbb{T}$ itself.
\end{remark}

In light of the definitions above, the level curves of $\gamma_{B^c}(x,y)$ are the
curves along which $e_{x+iy}(t,t_0)$ has a constant global exponential growth rate.
The level curves of $\omega_{B^c}(x,y)$ are the curves along which $e_{x+iy}(t,t_0)$
has a constant global (signal) frequency. As we will see, these two families of level
curves induce a natural coordinate system on $\mathbb{C}_{\mu(\mathbb{T})}$.

Using these definitions, we can reformulate $\mathcal{S}_{\mathbb{C}}(\mathbb{T})$ as
\begin{equation}\label{pswtype}
 \mathcal{S}_{\mathbb{C}}(\mathbb{T}) = \{x+iy \in \mathbb{C} \mid \gamma_{B^c}(x,y) <0\}.
\end{equation}
Equation \eqref{pswtype} is reminiscent of condition (C1) and $\mathcal{S}_\mathbb{C}(\mathbb{T})$
has been thought of as an average of Hilger circles \cite{DaGrMaJa}.

Finally, since $\gamma_{B^c}$ and $\omega_{B^c}$ are the average value of
$\myRe(\xi_{\mu(\tau)})$ and $\myIm(\xi_{\mu(\tau)})$, respectively, it makes sense to
study directly the complex function given by the average of $\xi_{\mu(\tau)}$.
In the next section, we will define and study the properties of a map $\overline{\xi}$,
called the \emph{ergodic cylinder transformation}, from the $xy$ plane to the
$\gamma \omega$ plane, i.e., a mapping from the pair $(x,y)$ to the pair
$(\gamma(x,y),\omega(x,y))$ for periodic time scales.


\section{The ergodic cylinder transformation}

We begin by defining the codomain of the ergodic cylinder transformation. We define
 \begin{gather*}
 \Omega := \sup_{x+iy \in B^c} \omega_{B^c}(x,y),\\
 \mathbb{C}^{\Omega} := \{z \in \mathbb{C} \mid -\Omega< \myIm(z) \leq \Omega\}.
 \end{gather*}
If $\Omega=\infty$, then $\mathbb{C}^{\Omega} = \mathbb{C}$.

\begin{definition}
\label{def:xibarBc} \rm
Assume \eqref{P}. Suppose $x+iy\in B^c$. We define the \emph{ergodic cylinder
 transformation off the branch cut}, $\overline{\xi}_{B^c}: B^c \to \mathbb{C}^{\Omega}$, by
$$
 \overline{\xi}_{B^c}(x+iy):=\frac{1}{L}\int_{0}^{L} \xi_{\mu(t)}(x+iy) \Delta t.
$$
\end{definition}

It is straightforward to show from \eqref{eq:exp_gamma_omega}, for $x+iy \in B^c$,
\begin{equation} \label{eq:gamOm}
 \overline{\xi}_{B^c}(x+iy) = \gamma_{B^c}(x,y) + i\,\omega_{B^c}(x,y).
\end{equation}

Our goal is to extend $\overline{\xi}_{B^c}$ to a map $\overline{\xi}:\mathbb{C}_{\mu(\mathbb{T})} \to \mathbb{C}^{\Omega}$
that induces an orthogonal curvilinear coordinate system on $\mathbb{C}_{\mu(\mathbb{T})}$ for which
the contours are given by the level curves of $\gamma$ and $\omega$. We will first show
that this is the case for $\overline{\xi}_{B^c}$, and then we will define the extension.

Analytic functions that are globally injective, or univalent, induce an orthogonal
coordinate system as the preimage of the rectangular coordinate system in the complex
plane via the Cauchy-Riemann equations. Therefore, we begin by showing
$\overline{\xi}_{B^c}:B^c \to \mathbb{C}^{\Omega}$ is analytic and univalent. To do so, we need the
following result, in the form first presented in \cite[Theorem~8]{AvAk}.

\begin{theorem}[Noshiro-Warschawski criterion]
If $f$ is analytic and nonconstant in a convex domain $D$, and
$$
\operatorname{Re}(e^{i \alpha} f'(z))>0
$$
for all $z \in D$ and for $\alpha \in \mathbb{R}$ fixed, then $f$ is univalent in $D$.
\end{theorem}

\begin{theorem}
\label{thm:analytic}
Suppose \eqref{P} holds. Then
\begin{enumerate}
 \item $\overline{\xi}_{B^c}:B^c \to \mathbb{C}^{\Omega}$ is analytic, and
 \item $\overline{\xi}_{B^c}:B^c \to \mathbb{C}^{\Omega}$ is univalent.
\end{enumerate}
\end{theorem}

\begin{proof}
Let $z=x+iy\in B^c$.

(1) Since the principal logarithm is analytic, $e_{z}$ is analytic for
$z \in \mathbb{C}_{\mu(\mathbb{T})}$ \cite{Ka}, and the nonregressive points lie in $B$,
it follows that $\overline{\xi}_{B^c}$ is analytic.

(2) To show univalence, from \cite{Ka}, the function $m_z:\mathbb{C}_{\mu(\mathbb{T})} \to \mathbb{C}$,
$$
m_z(s,t) = \int_{s}^{t} \frac{\Delta \tau}{1+z\mu(\tau)},
$$
is analytic in $z$, so
\begin{equation}
\begin{aligned}
\overline{\xi}_{B^c}'(z)
& = \frac{1}{L} \int_{0}^{L} \frac{\Delta \tau}{1+z \mu(\tau)}  \\
& = \frac{1}{L} \int_{0}^{L} \frac{(1+x \mu(\tau))\Delta \tau}{|1+(x+iy)\mu(\tau)|^2}
- \frac{i}{L} \int_{0}^{L} \frac{y \mu(\tau))\Delta \tau}{|1+(x+iy)\mu(\tau)|^2}.
\end{aligned} \label{calc}
\end{equation}

If $\mathbb{T}=\mathbb{R}$, univalence follows immediately since $\overline{\xi}_{B^c}(z) = z$.
If $\mathbb{T}\not=\mathbb{R}$, we consider three cases: $y<0,y=0$, and $y>0$. If $y>0$, let
$U := \{x+iy \in \mathbb{C} \mid y>0 \}$. Then, for $x+iy \in U$,
 \[
 \myRe(i\,\overline{\xi}_{B^c}'(x+iy))
  = \frac{1}{L} \int_{0}^{L} \frac{y \mu(\tau)\Delta \tau}{|1+(x+iy)\mu(\tau)|^2}>0,
 \]
since the integrand is nonnegative and not identically zero over $[0,L]_{\mathbb{T}}$ due to
the fact that $\mu(t) \not \equiv 0$, $y>0$, and (crucially) that
 $$
 \inf_{\tau \in [0,L]_{\mathbb{T}}} \frac{1}{|1+(x+iy) \mu(\tau)|^2} >0,
 $$
by \cite[Lemma 2.1]{Ka}. Since $U$ is convex, by the Noshiro-Warschawski criterion
(with $\alpha=\pi/2$), $\overline{\xi}_{B^c}$ is univalent on $U$.

Similarly, if $y<0$, $\overline{\xi}_{B^c}$ is univalent on the lower half-plane.
Note that if $z_1, z_2 \in \mathbb{C}$ with $\myIm(z_1)>0$ while $\myIm(z_2)<0$, then
$\overline{\xi}_{B^c}(z_1) \neq \overline{\xi}_{B^c}(z_2)$ since $\omega_{B^c}(x,y)>0$ for $y>0$
and $\omega_{B^c}(x,y)<0$ for $y<0$.

 Finally, if $y=0$ and $x> -1/\mu_{\max}$, $\omega_{B^c}(x,0) =0$ while
$\overline{\xi}_{B^c}'(x)>0$, which implies $\overline{\xi}_{B^c}$ is one-to-one for $x>-1/\mu_{\max}$.
Since $\omega_{B^c}$ is positive in the upper half-plane, negative in the lower-half
plane, and zero on the ray $(-1/\mu_{\max},\infty)$, and since $\overline{\xi}_{B^c}$ is
univalent on each of these spaces, it follows that $\overline{\xi}_{B^c}$ is globally
univalent.
\end{proof}

We now show that $\overline{\xi}_{B^c}$ can be extended to a map $\overline{\xi}:\mathbb{C}_{\mu(\mathbb{T})} \to \mathbb{C}$
which is globally univalent. We will call $\overline{\xi}$ the \emph{ergodic cylinder
transformation}. To do so, we treat the behavior along the branch cut $B$ carefully.
Since the branch cut consists of intervals between nonregressive points, and since
the properties of these intervals depend on properties of the time scale, we establish
some notation and technical lemmas to aid in the proof.

Let $M=\{\mu_n\}$ be the collection of distinct, positive graininesses of $\mathbb{T}\not=\mathbb{R}$,
ordered as
$$
\mu_{\max}=\mu_1>\mu_2>\cdots.
$$
Let $\mu_*=\inf_{n} \mu_n$ so that
$$
 -\frac{1}{\mu_*}<\cdots < -\frac{1}{\mu_2} < -\frac{1}{\mu_1}<0.
$$
Let
 \begin{gather*}
 I_n:=\Big(-\frac{1}{\mu_{n+1}},-\frac{1}{\mu_n}\Big),\ n=1,2,\dots,\\
 I_0:=
 \begin{cases}
 (-\infty,-1/\mu_*), & \mu_*>0,\\
 \emptyset, & \mu_*=0.
 \end{cases}
 \end{gather*}

\begin{lemma}
\label{lem:B_structure}
Assume \eqref{P}. Then $M$ is either finite or countably infinite. Moreover,
\begin{enumerate}
 \item If $M$ is finite,
 $$
 B\setminus \Big\{-\frac{1}{\mu_n}\Big\} = I_0 \cup \left(\bigcup_{n=1}^{N-1} I_n\right).
 $$

 \item If $M$ is countably infinite,
 $$
 B\setminus \Big\{-\frac{1}{\mu_n}\Big\} = \bigcup_n^\infty I_n.
 $$
\end{enumerate}
\end{lemma}

\begin{proof}
An arbitrary time scale has at most countably many graininesses.
 \begin{enumerate}
 \item If $M$ is finite, then $\mu_*>0$ so $I_0\not=\emptyset$ and the result follows.
 \item If $M$ is countably infinite, then we still have $ L=\sum_n \mu_n<\infty$,
 so it must be that $\mu_*=0$. Hence, $I_0=\emptyset$ and the result follows.
 \end{enumerate}
\end{proof}


\begin{theorem}
\label{thm:GlobalUnivalence}
$\overline{\xi}_{B^c}:B^c \to \mathbb{C}^{\Omega}$ can be extended to a globally univalent function
$\overline{\xi}:\mathbb{C}_{\mu(\mathbb{T})} \to \mathbb{C}^{\Omega}$.
\end{theorem}

\begin{proof}
We present the proof in two cases.
\smallskip

\noindent\textbf{Case 1: $M$ finite.}
Since $M$ is finite, we can list the $N$ distinct graininesses in order as
 \[
 \mu_{\max}=\mu_1>\mu_2>\mu_3>\cdots>\mu_N,
 \]
so that the corresponding non-regressive points satisfy
 \[
 -\frac{1}{\mu_N} < \cdots < -\frac{1}{\mu_2} < -\frac{1}{\mu_1} < 0.
 \]
Then, by Lemma \ref{lem:B_structure}, we have
$$
 B\setminus \Big\{-\frac{1}{\mu_n}\Big\}
 = I_0 \cup \left(\bigcup_{n=1}^{N-1} I_n\right).
$$
To extend $\overline{\xi}_{B^c}$, we consider
\begin{equation}\label{setup}
 \overline{\xi}_{B^c}^+(x):=\lim_{y\to 0^+}\overline{\xi}_{B^c}(x+iy),
 \quad\text{for } x\in B\setminus\{-1/\mu_n\}.
\end{equation}
Taking real and imaginary parts of $\overline{\xi}_{B^c}^+(x)$ gives
\begin{gather*}
 \gamma_B(x):=\operatorname{Re} \overline{\xi}_{B^c}^+(x)
 =\frac{1}{L}\int_{\{t\in[0,L]:\mu(t)=0\}} x\,\Delta t
  + \frac{1}{L}\int_{\{t\in[0,L]:\mu(t)>0\}} \frac{\ln|1+\mu(t) x|}{\mu(t)}
  \Delta t,
\\
\begin{aligned}
 \omega_B(x):=\operatorname{Im} \overline{\xi}_{B^c}^+(x)
 &=\frac{\pi}{L}\int_{\{t\in[0,L]:\mu(t)>0\}}
  \frac{\mathbf 1_{\{x<-1/\mu(t)\}}}{\mu(t)}\,\Delta t\\
 &=\frac{\pi}{L} \sum_{\{t\in[0,L]:\mu(t)>0\}} \mathbf 1_{\{x<-1/\mu(t)\}}\\
 &=\frac{\pi}{L}
 \begin{cases}
 \sum_{\{t\in[0,L]:\mu(t)\ge \mu_n\}} 1, & x\in I_n,\\
 \sum_{\{t\in[0,L]:\mu(t)>0\}} 1, & x\in I_0,
 \end{cases}
\end{aligned}
\end{gather*}
where $\mathbf{1}_A$ denotes the characteristic function on the set $A$.
Thus, $\omega_B$ is constant on each $I_n$; write $\omega_n:=\omega_B(x)$ for any
$x\in I_n$. Also, $\omega_B$ is constant on $I_0$ as well, so we let
$\omega_0:=\omega_B(x)$ for $x\in I_0$. Then
 $$
 \omega_1<\omega_2<\cdots<\omega_{N-1}<\omega_0.
 $$

For $x\in B\setminus\{-1/\mu_n\}$, differentiating under the integral yields
\begin{gather*}
\begin{aligned}
 \gamma_B'(x)
 &=\frac{1}{L}\int_{\{[0,L]:\mu(t)=0\}} 1\,\Delta t
 + \frac{1}{L}\int_{\{[0,L]:\mu(t)>0\}} \frac{1}{1+\mu(t) x}\,\Delta t\\
 &=p_0+\frac{1}{L}\int_{\{[0,L]:\mu(t)>0\}} \frac{1}{1+\mu(t) x}\,\Delta t,
\end{aligned}\\
\gamma_B''(x)
 =-\frac{1}{L}\int_{\{[0,L]:\mu(t)>0\}} \frac{\mu(t)}{(1+\mu(t) x)^2}\,\Delta t\ <\ 0,
\end{gather*}
where $$p_0:=\frac{1}{L}\int_{\{[0,L]:\mu(t)=0\}} 1\,\Delta t$$ is a constant.
Thus, $\gamma_B'$ is strictly decreasing on each $I_n$, $n=1,2,\dots$.
Moreover, on $I_n$, we have the end behavior
 \[
 \lim_{x\downarrow -1/\mu_{n+1}}\gamma_B'(x)=+\infty,\quad
 \lim_{x\uparrow -1/\mu_{n}}\gamma_B'(x)=-\infty.
 \]
Hence, there exists a unique $\alpha_n\in I_n$ with $\gamma_B'(\alpha_n)=0$.
Consequently,
 \[
 \gamma_B'(x)>0\ \text{on}\ \big(-\tfrac{1}{\mu_{n+1}},\,\alpha_n\big),
 \quad
 \gamma_B'(x)<0\ \text{on}\ \big(\alpha_n,\,-\tfrac{1}{\mu_n}\big).
 \]
We can similarly analyze the behavior of $\gamma_B$ on $I_0$. If $p_0=0$,
then $\gamma_B'(x)<0$ for $x\in I_0$. If $p_0>0$, then there exists a unique
$\alpha_0\in (-\infty,-1/\mu_N)$ such that $\gamma_B'(\alpha_0)=0$ and $\gamma_B'(x)>0$
for $x\in(-\infty,\alpha_0)$ while $\gamma_B'(x)<0$ for $x\in (\alpha_0,-1/\mu_N)$.

We define $\overline{\xi}_B:B\setminus\{-1/\mu_n\}\to\mathbb{C}$ by
 \begin{equation}\label{eff}
 \overline{\xi}_B(x)=
 \begin{cases}
 \gamma_B(x)+ i\,\omega_0, & p_0>0,\ x\in(-\infty,\alpha_0],\\[2pt]
 \gamma_B(x)- i\,\omega_0, & p_0>0,\ x\in(\alpha_0,-1/\mu_N),\\[2pt]
 \gamma_B(x)+ i\,\omega_0, & p_0=0,\ x\in(-\infty,-1/\mu_N),\\[2pt]
 \gamma_0(x)- i\,\omega_n, & x\in\big(-\tfrac{1}{\mu_{n+1}},\,\alpha_n\big],\;
  n=1,\dots,N-1,\\[2pt]
 \gamma_B(x)+ i\,\omega_n, & x\in\big(\alpha_n,\,-\tfrac{1}{\mu_n}\big),\ n=1,\dots,N-1,
 \end{cases}
 \end{equation}
Note that on each $I_n$, we use $-\overline{\xi}_{B^c}^+(x)$ for $x\in(-1/\mu_{n+1},\alpha_n]$,
but use $\overline{\xi}_{B^c}^+(x)$ for $x\in(\alpha_n,-1/\mu_n)$. As a preference, on $I_0$,
we use $\overline{\xi}_{B^c}^+(x)$ for $x\in(-\infty,\alpha_0]$, but use $-\overline{\xi}_{B^c}^+(x)$
for $x\in(\alpha_0,-1/\mu_N)$.

Finally, define $\overline{\xi}:\mathbb{C}_{\mu(\mathbb{T})} \to \mathbb{C}^{\Omega}$ by
 $$
 \overline{\xi}(z)=
 \begin{cases}
 \overline{\xi}_{B^c}(z), & z\in B^c,\\[2pt]
 \overline{\xi}_B(x), & z=x\in B\setminus\{-1/\mu_n\}.
 \end{cases}
 $$
We now show that $\overline{\xi}$ is globally univalent. We consider three subcases.
\smallskip

\noindent\textbf{Subcase 1: Two interior points.}
If $z_1,z_2\in B^c$ and $\overline{\xi}(z_1)=\overline{\xi}(z_2)$, then
$\overline{\xi}_{B^c}(z_1)=\overline{\xi}_{B^c}(z_2)$. By the univalence of $\overline{\xi}_{B^c}$ on
$B^c$, we get $z_1=z_2$.
\smallskip

\noindent\textbf{Subcase 2: One interior point, one branch cut point.}
Suppose $z\in B^c$ and $x\in B\setminus\{-1/\mu_n\}$.
The image $\overline{\xi}(B^c)=\overline{\xi}_{B^c}(B^c)$ is open.
The value $\overline{\xi}_B(x)$ is a one-sided (non-tangential) limit of $\overline{\xi}_{B^c}$
at $x\in\partial B^c$,
hence $\overline{\xi}_B(x)\in\partial\,\overline{\xi}_{B^c}(B^c)$ and is not an interior point.
Therefore $\overline{\xi}(z)\neq \overline{\xi}(x)$.
\smallskip

\noindent\textbf{Subcase 3: Two branch cut points.}
Let $x_1\neq x_2$ in $B\setminus\{-1/\mu_n\}$. Since $\overline{\xi}(x)=\overline{\xi}_B(x)$ on
$B\setminus\{-1/\mu_n\}$, it suffices to show $\overline{\xi}_B(x_1)\neq \overline{\xi}_B(x_2)$.

Suppose $x_1,x_2$ lie in one of the same half-intervals in \eqref{eff}.
Then, $\operatorname{Im} \overline{\xi}_B$ is constant while $\gamma_B'(x)$ has a fixed nonzero sign;
 hence $\overline{\xi}_B$ is strictly monotone along a horizontal line and
$\overline{\xi}_B(x_1)\neq \overline{\xi}_B(x_2)$.

On the other hand, if $x_1,x_2$ are in the two different halves of the same $I_n$,
then $\operatorname{Im} \overline{\xi}_B(x_1)=-\omega_n$ and $\operatorname{Im} \overline{\xi}_B(x_2)=+\omega_n$, so
$\overline{\xi}_B(x_1)\neq \overline{\xi}_B(x_2)$.

If $x_1,x_2$ lie in distinct intervals $I_j\neq I_n$, then
$\operatorname{Im} \overline{\xi}_B(x_\ell)\in\{\pm \omega_j,\pm \omega_n\}$ with $\omega_j\neq \omega_n$
since the sequence $\{\omega_n\}$ is strictly increasing in $n$. Again, the imaginary
parts differ, so the values are distinct.

Combining the three cases shows $\overline{\xi}$ is injective on
$\mathbb{C}_{\mu(\mathbb{T})} = \mathbb{C}\setminus\{-1/\mu_n\}$ when $M$ is finite.
\smallskip

\noindent\textbf{Case 2: $M$ countably infinite.}
If $M$ is countably infinite, by Lemma \ref{lem:B_structure},
 $$
 B\setminus \left\{-\frac{1} {\mu_n}\right\} = \bigcup_{n=1}^\infty I_n.
 $$
Since $I_0 = \emptyset$, we can define this simpler version of $\overline{\xi}_B:B\setminus\{-1/\mu_n\}\to\mathbb{C}$ by
\begin{equation}
 \overline{\xi}_B(x)=
 \begin{cases}
 \gamma_B(x)- i\,\omega_n, & x\in\big(-\tfrac{1}{\mu_{n+1}},\,\alpha_n\big],\\[2pt]
 \gamma_B(x)+ i\,\omega_n, & x\in\big(\alpha_n,\,-\tfrac{1}{\mu_n}\big),
 \end{cases}
\end{equation}
for all $n$ for which $\mu_{n+1}$ exists. Then
 \[
 \overline{\xi}(z)=
 \begin{cases}
 \overline{\xi}_{B^c}(z), & z\in B^c,\\[2pt]
 \overline{\xi}_B(x), & z=x\in B\setminus\{-1/\mu_n\}.
 \end{cases}
\]
The argument that $\overline{\xi}$ is injective in this situation is already contained in
the previous argument.
\end{proof}

\begin{thebibliography}{99}
\bibitem[1]{allouche1999ubiquitous} Jean-Paul Allouche, Jeffrey Shallit; \emph{The ubiquitous {P}rouhet-{T}hue-{M}orse sequence}, Sequences and their Applications: Proceedings of SETA'98, Springer, 1999, pp.~1--16.
\bibitem[2]{AvAk} FG~Avkhadiev, LA~Aksent'ev; \emph{The main results on sufficient conditions for an analytic function to be schlicht}, Russian Mathematical Surveys, \textbf{30} (1975), no.~4, 1.
\bibitem[3]{BoPeD} Martin Bohner, Allan Peterson; \emph{Dynamic {E}quations on {T}ime {S}cales: {A}n {I}ntroduction with {A}pplications}, Birkh\"auser, 2001.
\bibitem[4]{DaGrMaJa} John~M Davis, Ian~A Gravagne, Robert~J Marks, Billy~J Jackson; \emph{Regions of exponential stability for {LTI} systems on nonuniform discrete domains}, IEEE Proceedings of the 43rd Southeastern Symposium on System Theory, Auburn, AL (2011), 37--42.
\bibitem[5]{HiA} Stefan Hilger; \emph{Analysis on measure chains---a unified approach to continuous and discrete calculus}, Results in mathematics \textbf{18} (1990), no.~1, 18--56.
\bibitem[6]{JaDaE} Billy~J Jackson, John~M Davis, \emph{An ergodic approach to {P}\"otzsche-{S}iegmund-{W}irth's region of exponential stability}, Nonlinear Analysis: Hybrid Systems \textbf{36} (2020), 100848.
\bibitem[7]{JaDalap} Billy~J {J}ackson and John~M {D}avis; \emph{An ergodic approach to {L}aplace transforms on time scales}, Journal of Mathematical Analysis and Applications \textbf{502} (2021), no.~1, 125231.
\bibitem[8]{Ka} Ba{\c{s}}ak Karpuz; \emph{Analyticity of the complex time scale exponential}, Complex Analysis and Operator Theory, \textbf{11} (2017), 21--34.
\bibitem[9]{PoSiWi} Christian Potzsche, Stefan Siegmund, Fabian Wirth; \emph{A spectral characterization of exponential stability for linear time-invariant systems on time scales}, Discrete and Continuous Dynamical Systems, \textbf{9} (2003), no.~5, 1223--1242.
\bibitem[10]{PoDaGr} Dylan Poulsen, John Davis, Ian Gravagne; \emph{Mean square stability for systems on stochastically generated discrete time scales}, Nonlinear Analysis: Hybrid Systems, \textbf{31} (2019), 41--55.
\end{thebibliography}
\end{document}
